
% !TEX program = pdflatex
\documentclass[12pt,a4paper]{article}
\usepackage[a4paper,margin=1in]{geometry}
\usepackage[T1]{fontenc}
\usepackage[utf8]{inputenc}
\usepackage[expansion=false]{microtype}
\usepackage{setspace}
\usepackage{parskip}
\usepackage{titlesec}
\usepackage{xcolor}
\usepackage{listings}
\usepackage{fancyhdr}
\usepackage{lastpage}
\usepackage{enumitem}
\usepackage{hyperref}
\usepackage{graphicx}
\usepackage{float}

\onehalfspacing
\setlength{\parindent}{0pt}
\hypersetup{colorlinks=true,linkcolor=black,urlcolor=blue,
 pdftitle={Blockchain Laboratory - Experiment 5},pdfauthor={Nikhil Sunchu}}
\pagestyle{fancy}
\fancyhf{}
\lhead{Blockchain Laboratory}
\rhead{Experiment No. 5}
\cfoot{Page \thepage}
\setlength{\headheight}{15pt}
\titleformat{\section}{\large\bfseries}{\thesection.}{0.5em}{}
\titleformat{\subsection}{\normalsize\bfseries}{\thesubsection.}{0.5em}{}

\definecolor{codebg}{RGB}{248,248,248}
\definecolor{codecomment}{RGB}{80,120,80}
\definecolor{codekeyword}{RGB}{70,70,150}
\definecolor{codestring}{RGB}{150,70,70}
\definecolor{codetype}{RGB}{0,110,110}
\lstdefinestyle{soliditystyle}{
 language=C,
 backgroundcolor=\color{codebg},
 basicstyle=\ttfamily\footnotesize,
 keywordstyle=\color{codekeyword}\bfseries,
 commentstyle=\color{codecomment},
 stringstyle=\color{codestring},
 numbers=left, numberstyle=\tiny\color{gray}, numbersep=8pt,
 frame=single, framerule=0.4pt, breaklines=true,
 breakatwhitespace=false, showstringspaces=false,
 tabsize=4, columns=fullflexible, keepspaces=true,
 morekeywords={pragma,solidity,contract,mapping,address,uint,public,private,internal,external,view,pure,returns,return,function,constructor,if,else,for,while,require,msg,sender},
 keywordstyle=[2]\color{codetype}\bfseries,
 morekeywords=[2]{address,uint,bool,string,bytes}
}
\begin{document}
\begin{titlepage}
\centering
\vspace*{1.5cm}
{\Large\bfseries BLOCKCHAIN LABORATORY\par}
\vspace{1.2cm}
{\large\bfseries Experiment No. 5\par}
\vspace{0.8cm}
{\LARGE\bfseries Study of Solidity Program Structure, Basic Constructs, Compilation and Deployment Environment\par}
\vspace{1.5cm}
\begin{tabular}{rl}
\textbf{Name:} & Nikhil Sunchu \\[0.6cm]
\textbf{Roll No.:} & 31 \\[0.6cm]
\textbf{Class/section:} & D20A \\[0.6cm]
\textbf{Date:} & \hrulefill \\[0.6cm]
\end{tabular}
\vfill
{\large Blockchain Laboratory\par}
\end{titlepage}

\section{Aim}
Study the structure of a Solidity program, its basic constructs, and the compilation and deployment environment using Remix IDE.

\section{Theory}
\subsection{1. Solidity and Smart Contracts}
Solidity is a high-level, statically typed programming language used to develop smart contracts for Ethereum-compatible blockchain platforms. A smart contract is a program deployed on a blockchain that executes predefined logic when its functions are called. Remix IDE provides an online environment to write, compile, deploy, and test Solidity contracts.

\subsection{2. Structure of a Solidity Program}
A Solidity source file commonly contains an SPDX license identifier, a compiler-version pragma, a contract declaration, state variables, and functions. The \texttt{pragma} statement specifies the compiler versions compatible with the source code. A contract groups related state and behaviour under one name.

\subsection{3. Data Types, State Variables and Mappings}
Solidity provides types such as \texttt{uint} for unsigned integers, \texttt{bool} for Boolean values, \texttt{address} for Ethereum addresses, and \texttt{string} for text. State variables are stored as part of the contract's persistent state. A mapping stores key-value pairs and is useful for looking up a value using a key. In this experiment, a mapping associates an Ethereum address with an unsigned integer balance.

\subsection{4. Functions and Visibility}
Functions contain the operations performed by a contract. The \texttt{public} visibility specifier allows a function to be called both internally and externally. A function marked \texttt{view} can read contract state without modifying it. A function that changes a state variable is executed as a transaction when called externally.

\subsection{5. Compilation and Deployment}
Compilation checks the Solidity source code and generates contract artifacts, including the ABI and bytecode. Deployment creates a contract instance in the selected blockchain environment. Remix VM provides a simulated blockchain and test accounts, allowing contracts to be deployed and tested without using real Ether.

\subsection{6. Transactions and Gas}
A state-changing function call is submitted as a transaction and consumes gas for its execution. Remix displays information such as the transaction hash, block number, gas used, and execution cost. A read-only \texttt{view} function is generally invoked as a call and does not create a state-changing transaction.

\section{Tools and Environment Used}
\begin{itemize}
    \item Remix IDE (version 2.6.6 as shown in the screenshots)
    \item Solidity compiler, version 0.8.34 in the compiler panel
    \item Remix VM with the Osaka VM setting
    \item Solidity source file: \texttt{Nikhilsunchu.sol}
    \item Browser-based code editor, compiler, and Deploy \& Run Transactions panel
\end{itemize}

\section{Tasks Performed}
\subsection{Task 1 -- Write the Solidity Program}
\begin{itemize}
    \item Created a Solidity source file named \texttt{Nikhilsunchu.sol} in Remix IDE.
    \item Defined the \texttt{GasTransaction} contract.
    \item Used a mapping to store a balance for each Ethereum address.
    \item Implemented \texttt{setBalance()} to update the balance associated with the caller.
    \item Implemented \texttt{getBalance()} to read and return the caller's balance.
\end{itemize}

{\renewcommand{\baselinestretch}{0.9}\selectfont
\begin{lstlisting}[style=soliditystyle]
 // SPDX-License-Identifier: MIT
 pragma solidity ^0.8.18;

 contract GasTransaction {
     mapping(address => uint) public balances;

     function setBalance(uint amount) public {
         balances[msg.sender] = amount;
     }

     function getBalance() public view returns (uint) {
         return balances[msg.sender];
     }
 }
\end{lstlisting}
}

\subsection{Task 2 -- Compile the Contract}
\begin{itemize}
    \item Opened the Solidity Compiler panel in Remix IDE.
    \item Selected a compatible Solidity 0.8.x compiler.
    \item Compiled the source file by clicking the Compile button.
    \item Confirmed successful compilation before proceeding to deployment.
    \item The source code and compiler panel are shown in Figure~\ref{fig:compile}.
\end{itemize}

\begin{figure}[H]
\centering
\includegraphics[width=0.92\textwidth]{swappy-20261012-120105.png}
\caption{Solidity source code and compiler panel in Remix IDE}
\label{fig:compile}
\end{figure}

\subsection{Task 3 -- Deploy the Contract}
\begin{itemize}
    \item Opened the Deploy \& Run Transactions panel.
    \item Selected \texttt{Remix VM} as the execution environment.
    \item Selected and deployed the \texttt{GasTransaction} contract.
    \item Verified that the deployed contract appeared in the Deployed Contracts section.
    \item Used the displayed contract functions for interaction and testing.
\end{itemize}

\subsection{Task 4 -- Execute the State-Changing Function}
\begin{itemize}
    \item Entered \texttt{500} as the argument to \texttt{setBalance()}.
    \item Clicked the Transact button to execute the function.
    \item Remix created and mined a transaction because the function modifies contract state.
    \item Checked the transaction status, transaction hash, block hash, block number, and gas information.
    \item The successful transaction details are shown in Figure~\ref{fig:transaction}.
\end{itemize}

\begin{figure}[H]
\centering
\includegraphics[width=\textwidth]{swappy-20261012-120710.png}
\caption{Successful execution of setBalance(500) and transaction details}
\label{fig:transaction}
\end{figure}

\subsection{Task 5 -- Verify the Stored Value}
\begin{itemize}
    \item Called the \texttt{getBalance()} function from the deployed contract.
    \item Observed the returned value \texttt{500}.
    \item Confirmed that the state-changing transaction updated the mapping for the selected account.
    \item The returned value and Remix terminal output are shown in Figure~\ref{fig:balance}.
\end{itemize}

\begin{figure}[H]
\centering
\includegraphics[width=\textwidth]{swappy-20261012-120839.png}
\caption{Calling getBalance() and verifying the stored value}
\label{fig:balance}
\end{figure}

\section{Result}
The Solidity contract was compiled and deployed successfully in Remix VM. The \texttt{setBalance(500)} transaction updated the mapping for the caller's address, and \texttt{getBalance()} returned \texttt{500}. The experiment demonstrated the Solidity contract structure, state variables, mappings, function visibility, read-only functions, deployment, and transaction gas information.

\section{Conclusion}
The experiment provided practical familiarity with writing and compiling a Solidity smart contract in Remix IDE and deploying it in a simulated blockchain environment. By interacting with the \texttt{GasTransaction} contract, the operation of a mapping, a state-changing function, and a \texttt{view} function was observed. The transaction details also illustrated how Remix reports execution status and gas usage.
\end{document}
